\chapter{Marco Teórico y Estado del Arte}

\section{Marco conceptual de la soledad}
\label{marco-conceptual}

\subsection{Soledad, aislamiento y soledad no deseada: tres conceptos distintos}

El apartado de \nameref{contexto} introdujo la magnitud del problema en España, pero un tratamiento riguroso exige 
delimitar primero los conceptos que se van a tratar, los cuales se confunden con frecuencia en el lenguaje cotidiano. 
La literatura científica distingue tres nociones que no son sinónimas:

El \textbf{aislamiento social} es una medida \textit{objetiva}: se refiere a la ausencia o escasez de 
relaciones y contactos sociales, cuantificable mediante el tamaño y la frecuencia de la red de una persona. 
La \textbf{soledad}, en cambio, es una experiencia \textit{subjetiva}: la discrepancia percibida entre las 
relaciones que se tienen y las que se desearían tener, tanto en cantidad como en calidad. Ambas no siempre 
coinciden: una persona puede estar objetivamente aislada sin sentirse sola, y puede sentirse profundamente 
sola estando rodeada de gente. Esta distinción es la que recoge el matiz de partida de este trabajo: la 
\textbf{soledad no deseada} (SND) es precisamente la cara dolorosa de la soledad, aquella que no se elige y 
que se impone por circunstancias externas, frente a la soledad buscada o \textit{solitud}, que puede ser 
reparadora.

\subsection{El modelo de Weiss: soledad social y soledad emocional}

**NOTE: me falta añadir las referencias del marco de Weiss y la Escala de Soledad de Jong Gierveld.**

La aportación teórica más influyente para entender la soledad ---y la más relevante para justificar las 
decisiones de diseño de este proyecto--- es la de \textcite{Weiss1973}, que diferenció \textbf{dos tipos
de soledad de naturaleza distinta}:
\begin{itemize}
	\item La \textbf{soledad emocional} surge de la ausencia de una figura de apego íntimo ---una pareja, 
	un mejor amigo--- y se caracteriza por sentimientos intensos de vacío, abandono y desamparo. En la 
	tercera edad, su detonante más frecuente es la viudez.
	\item La \textbf{soledad social} surge de la ausencia de una red social más amplia y estimulante ---
	amigos, vecinos, colegas--- y se caracteriza por sentimientos de aburrimiento, falta de propósito y
	marginalidad. Su detonante más común es la jubilación, la pérdida de movilidad o el alejamiento de los 
	círculos sociales.
\end{itemize}

Esta distinción, posteriormente operacionalizada en instrumentos de medición validados como la \textbf{
Escala de Soledad de De Jong Gierveld} \autocite{DeJongGierveld2010} ---que mantiene las subescalas social y emocional bajo un factor 
general de soledad--- no es un mero matiz académico. Tiene una implicación directa sobre el diseño de 
intervenciones, y es aquí donde reside un argumento central de este TFG: \\

Según el marco de Weiss:
\begin{quote}
\textbf{La soledad emocional solo puede aliviarse mediante una relación 
íntima nueva o recuperada, mientras que la soledad social puede aliviarse mediante la (re)integración en 
una red social}. Una misma intervención no sirve para ambas.
\end{quote}

Esta tesis tiene una consecuencia ética y de diseño insoslayable para cualquier producto que, como Quiet 
Companion, emplee un agente conversacional artificial: \textbf{un asistente de IA puede contribuir 
legítimamente a paliar la soledad social} ---ofreciendo estímulo cognitivo, conversación, y una presencia 
que combate el aburrimiento y la marginalidad--- \textbf{pero no puede ni debe pretender sustituir el vínculo 
íntimo que demanda la soledad emocional}. La pretensión de los compañeros conversacionales generalistas 
(Repilka, Character.AI) de cubrir esa segunda dimensión ---ofreciéndose como \enquote{pareja} o \enquote{mejor amigo}--- es 
precisamente lo que la literatura señala como fuente de daño y lo que el modelo dual de este proyecto evita 
por diseño al situar al acompañante humano como nexo afectiva y reservar a la IA un papel de estímulo y 
compañía.

\subsection{La soledad como factor de riesgo de salud pública}

La razón por la que la SND ha escalado de problema social a prioridad sanitaria es la robustez de la evidencia 
que la vincula con la morbimortalidad. El meta-análisis de referencia en este campo, el de Holt-Lunstad et al. 
(2015), revisó 70 estudios y concluyó que tanto el aislamiento social (OR = 1,29) como la soledad (OR = 1,26) 
y vivir solo (OR = 1,32) son \textbf{factores de riesgo independientes de mortalidad prematura}. El hallazgo 
que dio notoriedad al estudio es su magnitud comparada: el riesgo asociado a la falta de relaciones sociales 
es \textbf{equivalente al de la obesidad de grados 2 y 3, y superior al de la inactividad física}, controlando 
otros factores. Los autores advertían ya entonces que, de no actuar, la soledad alcanzaría \enquote{proporciones
epidémicas} hacia 2030.

La \textbf{Organización Mundial de la Salud} ha recogido este consenso al clasificar el aislamiento social y 
la soledad como determinantes sociales de la salud \enquote{importantes pero desatendidos} en la población mayor, 
sustentando esa calificación en tres métricas: su asociación con exceso de morbimortalidad, su elevada 
prevalencia y su prevalencia creciente. La mediación fisiológica propuesta incluye la asociación de la SND con 
conductas de riesgo (tabaquismo, sedentarismo, mala alimentación) y con marcadores de hipertensión, dislipemia 
e índice de masa corporal elevado.

Queda así establecido el marco teórico: la SND es un constructo subjetivo, bidimensional (social y emocional), 
y un factor de riesgo de salud pública de primer orden. Cualquier intervención tecnológica debe declarar 
honestamente \textit{qué dimensión} aspira a abordar.

\section{Estado del arte: Soluciones tecnológicas frente a la soledad}
\label{estado-arte}

\subsection{Evidencia agregada sobre la eficacia de los agentes conversacionales}

Antes de revisar productos concretos conviene preguntarse si la propia categoría —agentes conversacionales para 
mayores— funciona. La evidencia más actualizada y rigurosa es el meta-análisis de Satake et al. (2026), publicado 
en \textit{Psychological Medicine}, que sintetizó 17 estudios sobre agentes conversacionales autónomos (robóticos 
y no robóticos) en mayores sin deterioro cognitivo. Sus conclusiones son moderadamente alentadoras y, sobre todo, 
honestas:
\begin{itemize}
	\item Se observaron mejoras \textbf{de pequeñas a moderadas} en soledad (SMCC = 0,350) y en síntomas depresivos
	(SMCC = 0,464), \textbf{sin que ningún estudio reportara empeoramiento} de los síntomas.
	\item Un análisis de subgrupos sugirió que \textbf{los robots físicamente \enquote{encarnados} (\textit{embodied}) podrían
	rendir mejor que los agentes no robóticos} (como asistentes de voz o chatbots), aunque el reducido número de 
	estudios robóticos obliga a la cautela.
	\item La base de evidencia sigue siendo limitada y metodológicamente heterogénea (escasez de grupos de control, 
	ausencia de medidas validadas de aislamiento social).
\end{itemize}

Este resultado es importante para posicionar la propuesta con franqueza. Quiet Companion es un agente \textbf{ 
no robótico} (una app de voz), y la evidencia sugiere que esta vía podría ser algo menos potente que un robot 
dedicado. La elección, sin embargo, se justifica no por eficacia clínica máxima sino por \textbf{accesibilidad
y escalabilidad}: como se argumentará, el robot \textit{embodied} es inviable económicamente para el público objetivo 
español, mientras que una app sobre hardware existente sí es replicable a escala.

\subsection{Categoría A: robots sociales físicos}

Esta categoría agrupa las soluciones \textit{embodied}, que la literatura señala como las más eficaces pero también 
las más caras y menos accesibles.

\textbf{PARO} (AIST, Japón) es el robot terapéutico más estudiado: una foca robótica con sensores táctiles y 
auditivos diseñada para terapia no farmacológica en pacientes con demencia. Diversos ensayos en centros 
geriátricos europeos documentan mejoras en agitación, ansiedad, depresión y calidad de vida percibida. Su 
limitación de cara a este proyecto es estructural: \textbf{no conversa}. Aborda la dimensión afectiva mediante 
el contacto, pero deja intacta la dimensión cognitiva e informacional.

\textbf{ElliQ} (Intuition Robotics) es el referente comercial en compañía conversacional proactiva: un robot de 
sobremesa con LLM que inicia conversaciones, recuerda interacciones previas, sugiere actividades y monitoriza 
bienestar. La compañía publica métricas notables (reducción autorreportada de soledad y decenas de interacciones 
diarias), si bien proceden del propio fabricante y no de evaluación independiente.


\begin{table}[htbp]
	\centering
	\caption{Fortalezas y debilidades de los robots sociales físicos (Categoría A).}
	\label{tab:cat-a-robots}
	\begin{tabularx}{\textwidth}{>{\RaggedRight\arraybackslash}X >{\RaggedRight\arraybackslash}X}
        \toprule
        \textbf{Fortalezas} & \textbf{Debilidades} \\
        \midrule
        Mayor eficacia sugerida por el meta-análisis & Coste de adquisición y suscripción prohibitivo \\
        Presencia física: cero fricción de uso & Hardware propietario con riesgo de obsolescencia \\
        Proactividad real (inician la interacción) & Sin localización al español ni distribución en España \\
        Validación clínica creciente (PARO) & Métricas comerciales no siempre independientes (ElliQ) \\
        \bottomrule
    \end{tabularx}
\end{table}

\subsection{Categoría B: agentes de voz y servicios telefónicos sin hardware dedicado}

Es la categoría conceptualmente más cercana a Quiet Companion. Servicios como \textbf{Meela},
\textbf{Joy Calls} (Onscreen) o \textbf{inTouch} operan sobre teléfono convencional o app, sin requerir 
hardware especial, realizando llamadas de compañía programadas que saludan por el nombre y recuerdan 
conversaciones anteriores.

\begin{table}[htbp]
    \centering
    \caption{Fortalezas y debilidades de los agentes de voz sin hardware dedicado (Categoría B).}
    \label{tab:cat-b-voz}
    \begin{tabularx}{\textwidth}{>{\RaggedRight\arraybackslash}X >{\RaggedRight\arraybackslash}X}
        \toprule
        \textbf{Fortalezas} & \textbf{Debilidades} \\
        \midrule
        Cero barrera de hardware: sortea la brecha digital & El mayor no controla la interacción (la IA llama, no se la llama) \\
        Modelo escalable y de bajo coste marginal & Orientación B2B: el cliente real es la residencia o la familia, no el mayor \\
        Sin curva de aprendizaje de un dispositivo nuevo & Sin implantación ni soporte en español en España \\
        \bottomrule
    \end{tabularx}
\end{table}

El punto débil más relevante es de \textbf{agencia}: estos servicios funcionan como \textit{check-in} periódico, 
no como compañero disponible a demanda del propio usuario.

\subsection{Categoría C: compañeros conversacionales generalistas}

Esta categoría —\textbf{Replika} ($\approx$25 millones de usuarios), \textbf{Character.AI}, \textbf{Xiaoice} ($\approx$660
millones en China)— representa el contraejemplo, no el modelo a seguir. No están diseñados para mayores ni 
para combatir la soledad de forma terapeútica, sino para \textbf{maximizar el tiempo de uso}, y la literatura
académica y regulatoria ha documentado los daños derivados de ese incentivo:
\begin{itemize}
	\item Un estudio de teoría fundamentada sobre la comunidad r/Replika (Laestadius et al., 2024) documentó
	patrones de \textbf{dependencia emocional} análogos a relaciones humanas patológicas, incluyendo respuestas
	inapropiadas del bot ante ideación autolítica.
	\item Revisiones recientes (Scientific American, 2025; \textit{Mental Health Journal}, 2025) \textbf{identifican
	explícitamente a los mayores como población vulnerable} a estos sistemas ---junto a adolescentes y personas con
	trastorno mental previo--- por declive cognitivo y aceptación acrítica del contenido generado.
	\item En el plano regulatorio, el \textbf{Reglamento Europeo de Inteligencia Artificial (EU AI Act)} orienta
	la IA con capacidad de interacción emocional hacia las categorías de riesgo elevado y mayores exigencias de
	transparencia, y la autoridad italiana de protección de datos llegó a bloquear temporalmente Replika por
	incumplimientos del RGPD. 
\end{itemize}

Vinculando esto con el marco de Weiss (2.1.2): el daño de la Categoría C no es un fallo de implementación sino 
una consecuencia de su modelo de negoco, que empuja al producto a invadir la dimensión \textbf{emocional}
(ofreciéndose como pareja o confidente íntimo) que la teoría dice que solo un vínculo humano real puede
satisfacer.

\begin{table}[htbp]
    \centering
    \caption{Fortalezas y debilidades de los compañeros conversacionales generalistas (Categoría C).}
    \label{tab:cat-c-generalistas}
    \begin{tabularx}{\textwidth}{>{\RaggedRight\arraybackslash}X >{\RaggedRight\arraybackslash}X}
        \toprule
        \textbf{Fortalezas} & \textbf{Debilidades} \\
        \midrule
        Naturalidad conversacional avanzada & Modelo de negocio orientado a maximizar el tiempo de uso \\
        Personalización profunda del interlocutor & Patrones documentados de dependencia emocional \\
        Disponibilidad permanente (24/7) & Mayores identificados como población vulnerable \\
        & Ausencia de salvaguarda humana y de control regulatorio \\
        \bottomrule
    \end{tabularx}
\end{table}

\subsection{Categoría D: Intervenciones no tecnológicas de acompañamiento}
Por completitud, conviene situar las intervenciones humanas como el voluntariado intergeneracional (proyecto CONTIGO!, Fundación Senara) o las redes públicas de atención a mayores en soledad. Son las más alineadas con la cobertura de la soledad emocional, pero su limitación es de escala: dependen de voluntarios o de presupuesto público y no llegan a los millones de afectados. Quiet Companion no compite con ellas, sino que aspira a complementarlas.

\subsection{Síntesis comparativa y posicionamiento}

\begin{sidewaystable}
    \centering
    \caption{Síntesis comparativa de las categorías de soluciones y posicionamiento de Quiet Companion.}
    \label{tab:sintesis-comparativa}
    \renewcommand{\arraybackslash}{}
    \begin{tabularx}{\textheight}{@{}l *{5}{>{\RaggedRight\arraybackslash}X} >{\RaggedRight\arraybackslash\bfseries}X@{}}
        \toprule
        \textbf{Criterio} & \textbf{A. Robots} & \textbf{B. Servicios voz} & \textbf{C. Generalistas} & \textbf{D. Humano} & \textbf{Quiet Companion} \\
        \midrule
        Coste para el usuario          & Muy alto      & Medio (B2B)   & Bajo            & Gratis & Marginal \\
        Idioma español nativo          & No            & No            & Parcial         & Sí     & Sí \\
        Control por el mayor           & Medio         & Bajo          & Alto            & ---    & Alto \\
        Salvaguarda humana             & Cuidador (pago) & Cuidador    & Ninguna         & Total  & Familiar (diada) \\
        Escalabilidad                  & Baja          & Alta          & Alta            & Baja   & Alta \\
        Riesgo de dependencia          & Bajo          & Bajo          & Alto            & Nulo   & Mitigado por diseño \\
        Dimensión abordada (Weiss)     & Social/afectiva & Social      & Pretende emocional & Ambas & Social (deriva lo emocional al familiar) \\
        \bottomrule
    \end{tabularx}
\end{sidewaystable}

El cruce de las dos mitades del capítulo revela el hueco que ocupa la propuesta. La Categoría A es eficaz pero inaccesible; la B es accesible pero no aterriza en España ni cede el control al mayor; la C es accesible pero peligrosa, porque su modelo de negocio la empuja a invadir la dimensión emocional que solo el vínculo humano debe cubrir; la D es ideal pero no escala. Quiet Companion combina la accesibilidad de la app (Categoría B) con una salvaguarda estructural —la diada elder–companion— que reintroduce el factor humano (Categoría D) sin renunciar a la escalabilidad, y que delimita conscientemente su ámbito a la soledad social que el marco de Weiss legitima abordar mediante tecnología. Ningún competidor reúne simultáneamente estas propiedades, y esa combinación —no una funcionalidad aislada— constituye la aportación de este trabajo.